\section{What the blind curve hid}
\label{sec:calibration}

Three acceptance criteria were fixed in a design document before any sweep ran:
a detectable halving of per-cell output, $s = 0.5$; interval coverage of 90\%;
and discrimination of 80\%, the probability the interval excludes zero at that
effect.\iffull{} The cutpoints follow mechanically --- the smallest $n$ meeting
both criteria, and the smallest meeting coverage alone.\fi The pre-registration is
mechanical rather than documentary: the panel, the labelling axes and the
falsification rule are frozen in a configuration file
\iffull
with a \texttt{frozen\_on} stamp and mirrored by an exact-value assertion in the
test suite, so continuous integration fails if anyone edits one without the
other.
\else
and asserted in the test suite, so continuous integration fails if anyone edits
one without the other.
\fi

Run properly on real held-out cells, the committed cutpoint of 50 does not come
back\iffull{} (Figure~\ref{fig:calibration})\fi. The generator filters eligibility on
patient alone, so both synthetic arms draw from the held-out patients' cells of
\emph{both} tissues pooled.
\iffull
The design document lists what stays fixed and never
mentions tissue. Nobody made that choice; an implementation detail made it, and
it changes the answer.
\else
The design document never mentions tissue: nobody made
that choice, an implementation detail made it, and it changes the answer.
\fi
So we report both readings rather than the one that
yields a cutpoint. Pooled, the routine returns no cutpoint at all: discrimination
never reaches the 0.80 target in any count bin, topping out at 0.78 across 13
seeds at 50 replicates and at 0.75 across 8 seeds at 500.\iffull{} Restricted to
the reference tissue it does return one, and that is where the second lesson of
this paper repeats itself.\else{} Restricted to the reference tissue it does
return one, and that is where this paper's second lesson repeats itself.\fi

\iffull
\begin{figure}[t]
\centering
\iffull
\includegraphics[width=\linewidth]{figures/fig1_calibration.pdf}
\else
\includegraphics[width=0.88\linewidth]{figures/fig1_calibration.pdf}
\fi
\iffull
\caption{\textbf{The committed cutpoint is not what the pre-registered procedure
returns, under either reading of an unspecified design choice.} Dashed lines mark
the targets fixed before the sweep. Shading marks the committed rule's three
states. \emph{Left:} with the draw pool our sweep code implements, the routine
returns no cutpoint at any tested count. \emph{Right:} under a second reading of
the same unspecified choice it does return one --- but the value is the grid's,
not the cohort's: $100$, $90$ and $70$ on three grids of increasing resolution,
against a committed 50 that none of them returns.}
\else
\caption{\textbf{The committed cutpoint is not what the pre-registered procedure
returns, under either reading of an unspecified design choice.} Dashed lines are
the targets fixed before the sweep; shading marks the committed rule's three
states; the committed value is 50 in both panels.}
\fi
\label{fig:calibration}
\end{figure}
\fi


\iffull\else
\paragraph{The grid was reporting its own resolution.} The committed grid varies
a mature \emph{fraction} against a fixed 2{,}000 cells, so it can reach no count
between 40 and 100, and returns $ok = 100$ on every seed. Adding nine fractions
inside that gap and changing nothing else moves the answer to 90, against a
committed 50: the crossing it was hiding runs $0.728$ at 45 mature cells,
$0.789$ at 65 and $0.858$ at 90. A grid that cannot resolve a crossing still
returns a stable number, which is \S\ref{sec:blind}'s failure in a second
statistic.
\fi
\iffull
\paragraph{The grid was reporting its own resolution.} Our committed grid varies
a mature \emph{fraction} against a fixed 2{,}000 cells, so the counts it can
reach are $\{0, 20, 40, 100, 200, 400, 800\}$ --- nothing between 40 and 100. On
the reference pool it returns $ok = 100$ on every one of the eight seeds at 500
replicates\iffull{}, and on twelve of thirteen at 50,\fi{} and that stability is
what made the number look like a measurement.

\iffull
The superseded internal report gave 120, and we cannot reproduce it: 120 is not
a count the committed grid can reach. Whatever produced it was not the committed
configuration, and the run is the uncommitted one of
the artifact's defect log item~3, so there is nothing left to re-derive it
from. We report that rather than reconcile it.
\fi

\textbf{And the reported cutpoint is not a count. It is a bin label.} Counts are
binned into twelve log-spaced intervals before coverage and discrimination are
computed, and the value reported for a bin is the \emph{median of the counts
that landed in it}. Over $20$--$800$ the edges are $\{0, 20, 28, 39, 55, 76,
107, 150, 209, 293, 409, 572, 800\}$, and they are identical for every grid
because they are derived from the same span. Nothing in the output says so, and
we did not know it while writing the paragraph above.

We added nine fractions inside the gap, reaching 50, 60, 70, 80, 90, 110, 120,
130 and 150, and changed nothing else. The answer moves to 90 --- and it moves
\emph{within one bin}. Both grids return the interval $(76, 107]$; the committed
grid held only $\{100\}$ there and the filled grid holds $\{80, 90, 100\}$, so
the median fell from 100 to 90 while the bin boundaries stayed put; the added counts change the evidence
aggregated within that bin. \textbf{At the resolution of the statistic that produced it, the answer did
not move at all.} What the filled grid did buy is a genuinely new bin at
$(55, 76]$, where discrimination reads $0.789$ --- under target, so the crossing
lies above it.

\textbf{A third grid, and the number keeps falling.} A denser grid reaching
below 20 --- attainable counts $\{2.5, 10, 17.5, 27.5, 42.5, 70, 105, 150, 400,
800\}$ --- returns $ok = 70$ on eight of eight seeds. Discrimination against
count is smooth, and the three grids agree on it wherever they overlap
(medians over eight seeds of $0.712$, $0.728$ and $0.726$ near 40, and $0.880$,
$0.858$ and $0.878$ near 100). What
differs is only which point each grid happens to carry first above the $0.80$
line:

\begin{center}\small
\begin{tabular}{@{}lccc@{}}
\toprule
 & \textbf{committed} & \textbf{filled} & \textbf{dense} \\
\midrule
crossing bracketed to & $(40, 100]$ & $(65, 90]$ & $(42.5, 70]$ \\
reported $ok$ & $100$ & $90$ & $\mathbf{70}$ \\
reported $wide$ & $40$ & $45$ & $42$ \\
\bottomrule
\end{tabular}
\end{center}

\textbf{The reported cutpoint decreases across these tested grids; this does not
establish monotonicity or non-convergence under further refinement. No tested
grid, cohort or seed returned the committed 50.} Each value is the first point above a crossing, reported as
though it were the crossing. Within any one grid $wide$ has \emph{exactly zero}
seed-to-seed variance --- one value on all eight seeds, in all six
cohort-by-grid cells we ran ($40$, $45$, $42$ here; $100$, $65$, $70$ on the
second cohort) --- so the quantity that looks most like a measurement is the one
whose stability comes entirely from having nowhere else to land. The pooled draw
returns no cutpoint under any of the three grids, on any seed of either cohort,
discrimination topping out at $0.750$ and $0.795$ against the $0.80$ target, so
the non-identifiability there is not a resolution artefact.

\textbf{We do not endorse a calibrated cutpoint: the procedure returns
grid-dependent candidates without patient-level uncertainty.} It returns its own grid, and it did so with a stability that
survived three grids, thirteen seeds and two cohorts without once warning us.
On the real cohort, rather than on the sweep, abstention is the modal outcome
under both population definitions at the resolution the compositional question
is defined on, and the coarse resolutions return \emph{not identifiable}
rather than a number.
\fi

\iffull
A second cohort agrees on the verdict and not the number: its pooled draw
returns no cutpoint on any seed, and its reference draw returns $400$ or
$800$, the top two counts its grid can express. Eight to sixteen times the
committed 50, by the same mechanism.

\paragraph{Three further statistics, and where they are.} The same failure
recurs three more times in this pipeline --- a confound diagnostic whose
tolerance sat above its own attainable bound, a pre-registered criterion that
could not pass when the hypothesis held, and a premise check whose sensitivity
was a property of the gene it controlled. Each survives correct code and each
is recorded in Appendix~\ref{app:guards} rather than argued here; none of this
paper's claims rests on them. The fourth is not a guard at all.

\paragraph{The same failure, one layer out: the interval.}
Every reading in this paper ends in a percentile bootstrap over patients. Conditional on the sample, the
bootstrap mean has standard deviation $s\sqrt{(n-1)/n}/\sqrt{n}$, with $s$
the sample standard deviation. A normal approximation gives the percentile interval
$\bar{x} \pm z\,s\sqrt{(n-1)/n}/\sqrt{n}$ against the normal-population Student-$t$ interval
$\bar{x} \pm t_{n-1}\,s/\sqrt{n}$, and the ratio of the two widths is
\begin{center}
$z\sqrt{(n-1)/n} \,/\, t_{n-1}$ \quad --- \quad \textbf{a function of $n$ alone,
with no data in it.}
\end{center}
Two errors running the same way: a normal quantile where a $t$ quantile is
needed, and the empirical-distribution variance in place of the sample-variance
estimate. At
$n = 44$ the interval is $0.96\times$ the width it claims and a nominal 5\% test
rejects 5.9\% of the time; at $n = 20$, $0.91\times$ and 7.1\%; at $n = 10$,
$0.82\times$ and 9.6\%; at $n = 4$, $0.53\times$ and 18.8\%. Simulation at 1{,}500 trials per cell agrees
with the closed form to within $2.7$ percentage points in every one of twenty
cells, against a Monte Carlo standard error of $0.6$ to $1.0$ points at these
rates. The simulated rate sits \emph{above} the closed form in 14 of the 20,
which is the direction the skew a normal approximation drops would produce, but
at this trial count that majority is what chance alone gives about one time in
seventeen --- so we report the closed form as an approximation the simulation
confirms, and not as a floor.

\textbf{Neither sophisticated repair works, and one is worse.} BCa is
miscalibrated in 17 of 20 cells against the percentile's 14, because its bias
correction and jackknife acceleration are estimated from the same ten numbers.
Raising the bootstrap replicate count does nothing at all: the error is not
Monte Carlo error, so more replicates estimate the wrong thing more precisely.
The plain Student-$t$ interval is calibrated in 19 of 20 cells and conservative
in the twentieth, worst case 5.9\%.

This is the same failure as the four above, one layer out. A guard that cannot
fail turns absence of evidence into a green light; an interval narrower than it
claims turns noise into a finding. Neither announces itself, and both are
functions of something other than the quantity a reader thinks they measure ---
\textbf{which is why the diagnostic this paper proposes is a question about the
statistic rather than a closer look at its output.} We did not find this by
inspecting intervals. We found it by asking what ours was a function of.

\fi
